\section{An exact change of sign in Suzuki's average over moduli}
\label{sec:fixed-suzuki}
Suzuki's equation (29), also printed in the \href{https://doi.org/10.1007/s11139-025-01289-y}{journal correction}, concerns
\[
 A_Q(x)=\sum_{q=3}^Q\left(
 \varphi(q)\sum_{\substack{n\le x\\n\equiv1\bmod q}}
       \frac{\Lambda(n)}{\sqrt n}\log\frac{x}{n}-4\sqrt x\right).
\]
His Theorem~7 gives negativity when $Q\ge x$. The paragraph after (30)
asks about a large fixed $Q$. The order of limits changes the answer.
Use the repository's Dirichlet GRH theorem throughout this section.

\paragraph{Theorem.}
For each fixed integer $3\le Q\le138$, $A_Q(x)$ is eventually negative.
For each fixed $139\le Q\le2{,}000{,}000$, it is eventually positive.
More generally, for any fixed $Q\ge139$ it is eventually negative if and
only if some primitive Dirichlet $L$-function of conductor at most $Q$
vanishes at $1/2$. In the other case it is eventually positive.
The starting value of $x$ may depend on $Q$.

\paragraph{The central residue.}
Set $Z_q(s)=\prod_{\chi\bmod q}L(s,\chi)$, $m_q=\operatorname{ord}_{1/2}Z_q$,
and define $C(q)$ by $Z_q'/Z_q(1/2+z)=m_q/z+C(q)+O_q(z)$.
Put $M(Q)=\sum_{q=3}^Q m_q$ and $B(Q)=\sum_{q=3}^Q C(q)$.
Character orthogonality and the twice-integrated explicit formula give
\[
 A_Q(x)=-\tfrac12 M(Q)(\log x)^2-B(Q)\log x+O_Q(1).
\]
In detail, the Mellin integrand for the $q$th inner weighted sum, including
its factor $\varphi(q)$, is $-(Z_q'/Z_q)(1/2+z)x^z/z^2$.
The pole at $z=1/2$ contributes $4\sqrt x$. The residue at zero yields
the two displayed powers and a constant. The noncentral critical zeros
yield a bounded oscillating sum because their denominators are
$|\Im\rho|^2$ and the modulus is fixed.
For primitive factors this is precisely the integrated explicit formula
used in Suzuki, \S\S3.1, 4.1 and 5.1
(\href{https://arxiv.org/html/2411.07436v3}{paper}).
For each Euler factor deleted on induction, expand its logarithmic
derivative in powers of $p^{-1/2-z}$. The coefficients at $z=0$ are
geometrically decreasing, so termwise summation with $\log(x/p^j)$
adds its stated linear term and a bounded remainder. This also justifies
the formula for the principal character.

\paragraph{An elementary coefficient valid at a central zero.}
Let $h=\varphi(q)$, $K=\mathbb Q(\zeta_q)$ and $D=|\operatorname{disc}K|$.
For $p^a\Vert q$ put $d=q/p^a$, $e=\varphi(p^a)$,
$f=\operatorname{ord}_d(p)$ with $f=1$ when $d=1$, and $g=h/(ef)$.
Euler products give $Z_q=\zeta_K\prod_{p\mid q}(1-p^{-fs})^g$.
The completed Dedekind zeta function, up to a fixed positive multiple,
is $D^{s/2}(2\pi)^{-hs/2}\Gamma(s)^{h/2}\zeta_K(s)$
and is symmetric about $1/2$. Its logarithmic derivative therefore has
zero constant Laurent coefficient at $1/2$, including at a zero.
The standard cyclotomic discriminant formula
$\log D=h(\log q-\sum_{p\mid q}\log p/(p-1))$
and $\Gamma'(1/2)/\Gamma(1/2)=-\gamma_{\rm E}-2\log2$ now give
\[
 C(q)=\frac{h}{2}\left(\log(8\pi)+\gamma_{\rm E}-\log q+
                         \sum_{p\mid q}\frac{\log p}{p-1}\right)
      +\sum_{p^a\Vert q}\frac{\varphi(q/p^a)\log p}{p^{f/2}-1}.
\]
Thus both the cutoff computation and the large-$q$ sign estimate apply
irrespective of uncomputed central values.

\paragraph{Certification of the crossover.}
The verifier \path{checks/suzuki_fixed_q.py}
encloses $B(Q)$ for every $3\le Q<10^4$. It finds
$45.431<B(138)<45.490$ and $-29.828<B(139)<-29.767$,
with $B(Q)>0$ throughout $3\le Q\le138$ and $B(Q)<0$ throughout
$139\le Q<10^4$. It uses rational remainder bounds for the logarithm
series after binary range reduction and integer square roots for the
ramified Euler factors. Machin's identity, the harmonic sum through 1024 with a trapezoid-rule
remainder, and the same logarithm series certify
$3.80138<\log(8\pi)+\gamma_{\rm E}<3.80140$.

For $q\ge10^4$, each $C(q)$ is negative. Indeed
\[
 \sum_{p\mid q}\frac{\log p}{p-1}
 \le\tfrac14\log q+\tfrac34\log2+\tfrac14\log3
 <\tfrac14\log q+0.8.
\]
Only the primes $2,3$ require the two corrections.
As $p^f\ge d+1$, we have
$p^{f/2}-1\ge(\sqrt2-1)\sqrt d$, also for $d=1$.
Since $e\ge p^a/2$ and $p^a\ge2$, the last sum in $C(q)$ divided by
$h$ is at most $(2+\sqrt2)\log q/\sqrt q$.
Consequently $C(q)/h<2.31-(3/8-0.035)\log q<0$ in this range.
It follows that $B(Q)<0$ for every $Q\ge139$.

The central orders are nonnegative. Imprimitive Euler factors do not
vanish at $1/2$, so $M(Q)>0$ exactly when a primitive character of
conductor at most $Q$ gives a central zero (the zeta value at $1/2$ is
nonzero). Platt's independent computation (Theorem 9.2 in the author manuscript)
verifies central nonvanishing
through modulus $2{,}000{,}000$
(\href{https://doi.org/10.1090/mcom/3077}{Math. Comp. 85 (2016), 3009--3027}).
Its central-value bound is distinct from the finite-height GRH bound in
the same paper. Substitution in the fixed-$Q$ expansion proves the theorem.

\paragraph{The second cutoff outside a set of finite logarithmic measure.}
Suzuki's equation (30) uses the sum
\[
 U_Q(x)=\sum_{q=3}^Q\varphi(q)
 \sum_{\substack{n\le xe^2\\n\equiv1\bmod q}}
 \frac{\Lambda(n)}{\sqrt n}\log\frac{x}{n}.
\]
For every fixed $Q$ its sign agrees with the eventual sign of $A_Q$ outside
a set $\mathcal E_Q\subset[2,\infty)$ with
$\int_{\mathcal E_Q}dx/x<\infty$.
To see this, put $y=xe^2$ and
$D_Q(y)=\sum_{q=3}^Q\varphi(q)\sum_{n\le y,\,n\equiv1\bmod q}
\Lambda(n)/\sqrt n$. The unsmoothed explicit formulas (36), (43), (44)
give
\[
 D_Q(e^t)=2(Q-2)e^{t/2}-M(Q)t+R_Q(t),
 \qquad\sup_{b\ge1}\int_b^{b+1}|R_Q(t)|^2\,dt<\infty.
\]
Indeed each noncentral zero contributes a frequency with coefficient of
absolute value $1/|\Im\rho|$ per unit of multiplicity. The fixed-conductor
zero count in a unit block is $O_Q(\log(2+|k|))$. The Fourier cutoff estimate
for unit intervals in Section~\ref{sec:suzuki-density} thus applies to the
finite family. The finitely many deleted Euler factors contribute bounded
terms through geometrically decreasing prime-power sums. Endpoints at
prime powers form a countable exceptional set.
From $U_Q(x)=A_Q(y)+4(Q-2)\sqrt y-2D_Q(y)$ and $\log y=\log x+2$
we obtain, writing $u=\log x$,
\[
 U_Q(e^u)=-\tfrac12M(Q)u^2-B(Q)u+H_Q(u),
 \qquad \sup_{b\ge1}\int_b^{b+1}|H_Q(u)|^2\,du<\infty .
\]
Since $B(Q)\ne0$, the absolute value of the polynomial part exceeds
$c_Q k$ on $[k,k+1]$ for all sufficiently large integers $k$, for some $c_Q>0$.
The set of sign failures in this interval has measure $O_Q(k^{-2})$.
Summing proves the finite logarithmic measure assertion.

\paragraph{A growing range.}
There is also a direct answer for part of Suzuki's growing-$Q$ question.
The prime weights in (29) are nonnegative and
$\sum_{q\mid n-1}\varphi(q)=n-1$ for $n\ge2$. Thus, uniformly in integers $Q\ge3$,
\[
 A_Q(x)\le \sum_{n\le x}\Lambda(n)\sqrt n\log(x/n)-4(Q-2)\sqrt x
 = (4/9+o(1))x^{3/2}-4(Q-2)\sqrt x.
\]
Partial summation with the prime number theorem gives the last equality,
using $\int_0^1 u^{1/2}\log(1/u)\,du=4/9$.
For every fixed $\varepsilon>0$ this proves eventual negativity throughout
$Q\ge(1/9+\varepsilon)x$, directly from the positive-sum argument of
Suzuki's proof. The result holds independently of RH.

\paragraph{What remains.}
Eventual negativity for arbitrarily large \emph{fixed} $Q$ in (29)
is equivalent to the existence of a single central Dirichlet zero.
This identifies the precise obstruction for that question and gives the
opposite sign over the verified conductor range starting at $139$.
The transition in intermediate growing ranges, and eventual pointwise sign for
(30), require additional arguments.
